% =====================================================================
%  Вычисление аналитической метрики на большом графе (out-of-core LeaderRank)
%  Итоговый отчёт по вступительному заданию.
%  Engine: pdfLaTeX (TeX Live). Кириллица: T2A + inputenc(utf8) + babel.
%  Сборка: ./build.sh   (pdflatex прогоняется дважды ради оглавления и ссылок)
% =====================================================================
\documentclass[12pt,a4paper]{article}

% --- Кодировка и язык (pdfLaTeX + кириллица) ---
\usepackage[T2A]{fontenc}
\usepackage[utf8]{inputenc}
\usepackage[english,russian]{babel} % основной язык — русский (последний в списке)

% --- Математика ---
\usepackage{amsmath}
\usepackage{amssymb}

% --- Вёрстка и оформление ---
\usepackage[a4paper,margin=2.5cm]{geometry}
% expansion=false: font expansion требует особой настройки шрифтов; оставляем
% только protrusion (визуальное выравнивание полей).
\usepackage[protrusion=true,expansion=false]{microtype}
\usepackage{enumitem}
\usepackage{booktabs}

% Немного свободы переносу строк, чтобы плотные формулы и моноширинные имена
% не вылезали на поля.
\setlength{\emergencystretch}{3em}

% --- Ссылки (загружается одним из последних) ---
\usepackage[hidelinks]{hyperref}

% --- Удобные сокращения для обозначений ---
\newcommand{\In}{\operatorname{In}}        % множество входящих соседей
\newcommand{\Out}{\operatorname{Out}}       % множество исходящих соседей
\newcommand{\degout}{\deg^{\mathrm{out}}}   % исходящая степень
\newcommand{\degin}{\deg^{\mathrm{in}}}     % входящая степень
\newcommand{\one}{\mathbf{1}}               % вектор из единиц
\newcommand{\T}{^{\top}}                     % транспонирование
\newcommand{\bigO}{\mathcal{O}}             % O-большое

\title{Вычисление аналитической метрики на большом графе\\[2pt]
       \large Out-of-core LeaderRank с параллельной обработкой}
\author{} % укажите ФИО автора
\date{\today}

\begin{document}
\maketitle
\begin{abstract}
\noindent
В работе вычисляется метрика центральности \textbf{LeaderRank} на ориентированном невзвешенном
графе, который не помещается в оперативную память. Рёбра хранятся на диске и читаются потоком, а в
памяти остаются только массивы размером с число вершин. Счёт идёт во много потоков и не зависит от
их числа: результат получается один и тот же бит в бит. В отчёте разобраны выбор метрики, варианты
хранения графа, устройство решения, анализ сложности и точности, а также замеры на синтетических и
реальных графах.
\end{abstract}
\medskip

\section{Введение}
\label{sec:intro}

\subsection{Задача}
Нужно посчитать аналитическую метрику на ориентированном невзвешенном графе. Главная сложность в
том, что граф целиком не влезает в оперативную память. В условии приведён пример: памяти $128$~МБ, а
на граф нужно около $1$~ГБ. Значит, держать весь граф в ОЗУ нельзя, надо работать с диском.

Условие добавляет ещё несколько требований.
\begin{itemize}[nosep]
  \item \textbf{Обработка с диска (out-of-core).} В памяти можно держать только данные порядка числа
        вершин, но не порядка числа рёбер.
  \item \textbf{Гипер-узлы.} Встречаются вершины со степенью на порядки выше средней (средняя степень
        $50$, а у одной вершины $50\,000$). Решение не должно ломаться на таких вершинах.
  \item \textbf{Параллелизм.} Узел один, но с несколькими ядрами. Нужно занять все ядра.
  \item \textbf{Достоверность.} Результат не должен зависеть от случайности и должен повторяться от
        запуска к запуску.
  \item \textbf{Разумное время.} На больших графах счёт должен укладываться в приемлемое время.
\end{itemize}

\paragraph{Формат ввода и вывода.} Он зафиксирован условием. Вход --- CSV с заголовком, первые два
столбца \texttt{from,to}, идентификаторы вершин типа \texttt{int32}, ребро направлено
\texttt{from}~$\to$~\texttt{to}, граф невзвешенный. Выход --- CSV со столбцами \texttt{vertex,rank},
по одному числу на вершину.

\subsection{Что сделано}
Выбрана метрика \textbf{LeaderRank} --- вариант PageRank без настраиваемых параметров. Решение
написано на Java (OpenJDK~21). Рёбра графа лежат на диске и читаются потоком, в памяти живут только
массивы размера $\bigO(n)$. Счёт распараллелен по ядрам и даёт один и тот же результат при любом
числе потоков. Лимит памяти задаётся снаружи (через \texttt{-Xmx} и cgroup), а программа сама
подстраивается под него. Подробные инструкции по сборке и запуску вынесены в \texttt{README},
здесь же описано устройство решения.

\subsection{Как читать отчёт}
Раздел~\ref{sec:metric} --- про метрику LeaderRank и почему выбрана именно она.
Раздел~\ref{sec:approaches} --- варианты хранения графа и выбор одного из них.
Разделы~\ref{sec:arch}--\ref{sec:memory} --- устройство решения: архитектура, параллелизм, память.
Раздел~\ref{sec:complexity} --- анализ сложности и точности.
Раздел~\ref{sec:testing} --- тесты и замеры.
Раздел~\ref{sec:conclusion} --- выводы.

\section{Метрика LeaderRank}
\label{sec:metric}

\subsection{Что это и как работает}
LeaderRank --- метрика центральности, то есть она даёт каждой вершине одно число
<<важности>>. Это модификация PageRank, придуманная для поиска лидеров в социальных
сетях. Главное отличие от PageRank в том, что у неё нет настраиваемых параметров.

Идея простая. В граф добавляют одну искусственную вершину --- \emph{фоновый узел} $g$ (ground node).
Его соединяют двусторонними рёбрами со всеми остальными вершинами ($g\to i$ и $i\to g$ для каждой
вершины $i$). По такому расширенному графу из $n+1$ вершины запускают обычное случайное блуждание,
без телепортации и без коэффициента затухания.

Почему это работает. Фоновый узел связан со всеми в обе стороны, поэтому расширенный граф сильно
связный и апериодичный. Из теории цепей Маркова отсюда следует, что у блуждания есть единственное
стационарное распределение, а степенной метод к нему сходится. И всё это без подбора параметров.
Заодно фоновый узел играет ту же роль, что телепортация в PageRank, но встроен прямо в структуру
графа.

\subsection{Формула}
Пусть у вершины $i$ исходящая степень $\degout(i)$. В расширенном графе добавляется ребро к фоновому
узлу, поэтому исходящая степень становится $\degout(i)+1$. На каждом шаге каждая вершина $i$ отдаёт
по $s_i/(\degout(i)+1)$ каждому реальному соседу и столько же фоновому узлу, а фоновый узел отдаёт по
$s_g/n$ каждой из $n$ вершин. Начальное состояние: $s_i=1$ у всех реальных вершин, $s_g=0$. Суммарная
<<масса>> равна $n$ и сохраняется на всех шагах. После сходимости масса фонового узла поровну
возвращается вершинам: $s_i \mathrel{+}= s_g/n$.

Удобно ввести вклад вершины, который считается раз за итерацию:
\[
  \mathrm{contrib}[j] = \frac{s_{\mathrm{old}}[j]}{\degout(j)+1}.
\]
Тогда обновление состояния записывается так:
\begin{align}
  \mathrm{base} &= s_{\mathrm{old}}[g]/n, \label{eq:base}\\
  s_{\mathrm{new}}[i] &= \mathrm{base} + \sum_{j\,:\,j\to i}\mathrm{contrib}[j],
      \qquad i=1,\dots,n, \label{eq:pull}\\
  s_{\mathrm{new}}[g] &= \sum_{j=1}^{n}\mathrm{contrib}[j]. \label{eq:ground}
\end{align}
Формулы~\eqref{eq:base} и~\eqref{eq:ground} --- дешёвые операции порядка $\bigO(n)$. Дорогая часть
одна --- сумма~\eqref{eq:pull} по рёбрам для каждого приёмника. В матричном виде это
$s_{\mathrm{new}} = \mathrm{base}\cdot\one + A\T\cdot\mathrm{contrib}$, то есть умножение
транспонированной матрицы смежности на вектор. Важно, что $2n$ рёбер фонового узла физически не
хранятся: его вклад считается двумя формулами.

\subsection{Итерация --- это один проход по рёбрам}
Ключевой факт: умножение разреженной матрицы на вектор --- это один проход по списку рёбер. Чтобы
посчитать $A\T\cdot\mathrm{contrib}$, достаточно для каждого ребра $j\to i$ прибавить
$\mathrm{contrib}[j]$ к ячейке $i$. Матрицу в памяти держать не нужно, нужен только поток рёбер.
Именно на этом держится всё решение: рёбра можно читать с диска по очереди, а в памяти хранить лишь
векторы размера $\bigO(n)$.

Останавливаемся, когда состояние почти перестаёт меняться. Критерий --- норма $L_1$ разности
$\sum_i |s_{\mathrm{new}}[i] - s_{\mathrm{old}}[i]| < \varepsilon$. Про конкретный порог и число
итераций --- в разделе~\ref{sec:complexity}. Масса сохраняется, поэтому нормировать вектор на каждом
шаге не нужно. Редкая перенормировка к сумме $n$ применяется только чтобы не накапливалась ошибка
округления.

\subsection{Где это полезно}
LeaderRank ищет влиятельные узлы. Исходно его применяли к социальной сети Delicious для поиска
лидеров мнений. По сути он годится везде, где применим PageRank: ранжирование
страниц, важность пользователей, графы цитирования, рекомендации. Его берут, когда не хочется
подбирать коэффициент затухания или когда в графе много висячих вершин.

\subsection{Почему выбрана именно она}
\label{sec:why-leaderrank}
Условие разрешает любую метрику, применяемую ко всему графу, и рекомендует PageRank, LeaderRank,
Katz, Jaccard, Dice. LeaderRank выбран по нескольким причинам.
\begin{enumerate}[itemsep=3pt]
  \item \textbf{Подходит формат.} LeaderRank даёт одно число на вершину, что прямо ложится в вывод
        \texttt{vertex,rank}. Пример выходного файла в условии так и называется ---
        \texttt{leader\_rank.csv}.
  \item \textbf{Естественно считается с диска.} Каждая итерация --- один последовательный проход по
        рёбрам. Значит рёбра (это и есть тот <<гигабайт>>) читаются потоком с диска, а в памяти
        остаются только векторы $\bigO(n)$. Память ограничена числом вершин, а не числом рёбер.
  \item \textbf{Не боится гипер-узлов.} Списки смежности не материализуются в памяти. Гипер-узел
        степени $50\,000$ --- это просто $50\,000$ записей в потоке рёбер на диске. Вклад вершины
        делится на её степень, никакого взрыва нет.
  \item \textbf{Легко параллелится.} Проход по рёбрам просто делится между потоками.
  \item \textbf{Детерминирован.} Случайности в алгоритме нет. При фиксированном порядке сложения
        результат повторяется.
  \item \textbf{Выигрывает у соседей по семейству.} В отличие от PageRank, не нужен параметр $d$ и
        не нужна отдельная обработка висячих вершин --- её берёт на себя фоновый узел. В отличие от
        Katz, не нужен спектральный радиус. Меры Jaccard и Dice вообще не подходят: требуют пересечения списков смежности (случайный доступ --- худший случай
        для диска) и катастрофически взрываются на гипер-узлах (узел степени $D$ участвует в
        $\bigO(D^2)$ парах). При этом схема вычислений у LeaderRank та же, что у PageRank, поэтому при
        желании PageRank поддерживается почти тем же кодом.
\end{enumerate}

\section{Способы представления графа}
\label{sec:approaches}

\subsection{Рамки решения}
\label{sec:scope}
Прежде чем выбирать представление, стоит зафиксировать рамки, в которых решается задача. Рёбер
$\bigO(m)$ --- это и есть <<гигабайт>> из условия, они не обязаны помещаться в память и живут на
диске. А структуры размером с число вершин, $\bigO(n)$ (векторы значений, степени, указатели),
наоборот, обязаны помещаться в ОЗУ: без них повершинную метрику не посчитать. Граница такая ---
$\bigO(m)$ не ограничено, $\bigO(n)$ обязано влезать.

Отсюда сразу отпадают подходы, которые выгружают на диск даже сами векторы значений $\bigO(n)$,
например Parallel Sliding Windows из GraphChi. Для этой задачи они избыточны: при <<гигабайте>> рёбер
вершин порядка $2$--$3$ млн, а их векторы --- это лишь десятки МБ, и они спокойно помещаются в память.
Случай, когда не влезает даже $\bigO(n)$, --- это уже другая задача, и она вне рамок данной работы.

\subsection{Push или pull}
Ядро метрики --- сумма по рёбрам~\eqref{eq:pull}. Организовать этот проход можно двумя способами, и
именно способ обработки определяет, как выгодно хранить граф.

\paragraph{Push (по источнику).} Обход идёт по рёбрам, сгруппированным по источнику $j$, и
$\mathrm{contrib}[j]$ прибавляется в ячейку приёмника $s_{\mathrm{new}}[i]$. Хранение здесь простое:
рёбра можно держать как есть, в порядке поступления, сортировать ничего не надо. Но при параллельном
проходе несколько потоков добавляют в одну и ту же ячейку популярного приёмника одновременно ---
возникает гонка. Пришлось бы ставить блокировки или атомарные операции, а порядок атомарных сложений
не определён, поэтому результат менялся бы в младших разрядах от запуска к запуску.

\paragraph{Pull (по приёмнику).} Обход идёт по приёмникам $i$, и для каждого собирается сумма
$s_{\mathrm{new}}[i] = \mathrm{base} + \sum_{j\to i}\mathrm{contrib}[j]$. У каждой ячейки результата
ровно один писатель, гонок нет. Зато появляется цена на хранении: чтобы входящие рёбра приёмника $i$
лежали подряд, рёбра нужно один раз сгруппировать и отсортировать по приёмнику. Это разовая работа на
этапе предобработки.

\subsection{Выбор}
Выбран pull. Решающим стало требование воспроизводимости. Push не требует сортировки, но его гонки и
неопределённый порядок сложения несовместимы с условием <<результат не должен зависеть от
случайности>>. Pull стоит разовой сортировки по приёмнику, зато на каждой из десятков итераций даёт
проход без блокировок, последовательное чтение с диска и детерминированный результат. Разовая
сортировка окупается уже на первых итерациях.

Из выбора pull прямо следует формат хранения: граф группируется по приёмнику и сортируется. Как он
раскладывается на диске и как строится за несколько внешних проходов --- в разделе~\ref{sec:arch}.

\section{Архитектура решения}
\label{sec:arch}

\subsection{Два этапа}
Решение состоит из предобработки (один раз) и итераций (много раз).
\begin{itemize}[nosep]
  \item \textbf{Предобработка} переводит входной CSV в компактный бинарный вид, удобный для
        потокового чтения. Это разовая плата, которая окупается: бинарный файл меньше CSV и упорядочен
        так, что итерация сводится к последовательному чтению.
  \item \textbf{Итерация} --- это шаг степенного метода LeaderRank~\eqref{eq:pull}. Бинарный файл
        читается потоком, считается ядро. В памяти живут только векторы $\bigO(n)$.
\end{itemize}
После сходимости плотные идентификаторы переводятся обратно в исходные \texttt{int32}, и результат
пишется в выходной CSV.

\subsection{Формат хранения}
\label{sec:format}
Ядру нужно для каждого приёмника $i$ перечислить его входящих соседей $j$. То есть нужна
транспонированная матрица $A\T$ --- граф, сгруппированный по приёмнику. Храним его как разреженный
столбцовый формат (CSC), разделённый на две части по принципу <<что влезает в ОЗУ --- в ОЗУ, что нет
--- на диск>>.
\begin{itemize}[itemsep=2pt]
  \item \texttt{sourcesPtr} --- массив \texttt{int[n+1]} указателей. Входящие рёбра приёмника $i$
        лежат в файле от позиции \texttt{sourcesPtr}$[i]$ до \texttt{sourcesPtr}$[i{+}1]$. Размер
        $\bigO(n)$, в ОЗУ.
  \item \texttt{sources} --- поток источников $j$, сгруппированных по приёмнику. Размер $\bigO(m)$,
        это и есть <<гигабайт>>, лежит на диске и читается потоком. Каждое ребро --- это $4$-байтовое
        целое (little-endian), только источник. Приёмник восстанавливается по положению в файле через
        \texttt{sourcesPtr}, поэтому повторно он не хранится.
  \item \texttt{outdeg} --- массив \texttt{int[n]} исходящих степеней (нужен для знаменателя
        $\degout(j)+1$). В ОЗУ.
  \item \texttt{originalIds} --- соответствие <<плотный id $\to$ исходный \texttt{int32}>> для вывода.
        В ОЗУ, нужен только при записи результата.
\end{itemize}

\subsection{Конвейер предобработки}
\label{sec:preprocess}
Предобработка --- это по сути внешняя сортировка: рёбра перегруппировываются по приёмнику, не
загружая их все в память. Именно эту <<работу с памятью>> условие требует написать самим. Этапы:
\begin{enumerate}[itemsep=3pt]
  \item \textbf{Один проход по CSV: переименование id и подсчёт степеней.} Исходные \texttt{int32}
        разрежены, поэтому они переводятся в плотный $0..n{-}1$ хеш-таблицей
        <<\texttt{int}$\to$\texttt{int}>> на примитивах. Таблица размера $\bigO(n)$ (вершины влезают в
        память, в отличие от рёбер). Заодно считаются входящие степени (для \texttt{sourcesPtr}) и
        исходящие (для знаменателя). Порядок
        присвоения плотных id фиксирован, по первому появлению. Это часть детерминизма. Чтобы не
        разбирать CSV второй раз, на этом же проходе рёбра в плотных id пишутся во временный бинарный
        файл.
  \item \textbf{Планирование корзин.} Диапазон приёмников делят на корзины так, чтобы в каждую
        попадало не больше $M$ рёбер. Размер $M$ подбирается под доступную память. Если у одной
        вершины входящая степень больше $M$, она не делится и попадает в отдельную <<большую>> корзину.
  \item \textbf{Раздача рёбер по корзинам.} Временный бинарный файл читается, каждое ребро
        дописывается в файл своей корзины. Записи последовательные, буферы фиксированного размера.
        Если корзин больше, чем можно открыть файлов под лимитом памяти, раздача идёт волнами:
        источник перечитывается несколько раз, по группе корзин за проход.
  \item \textbf{Сортировка корзин и сборка файла \texttt{sources}.} Обычная корзина целиком
        загружается в память, сортируется по ключу <<(приёмник, источник)>> и дописывается в итоговый
        файл. Большая корзина (гипер-узел) сортируется внешней сортировкой слиянием: режется на куски
        по $M$, каждый кусок сортируется и сливается потоком. Так память на одну корзину остаётся
        $\bigO(M)$ при любой степени вершины.
\end{enumerate}

Итог предобработки --- файл \texttt{sources} на диске плюс массивы $\bigO(n)$ в памяти. Дальше
итерация работает только с ними.

\paragraph{Почему детерминированно.} Итоговая раскладка (по приёмнику, источники по возрастанию) ---
это чистая функция от графа. Она не зависит ни от $M$, ни от границ корзин, ни от числа потоков. Это
проверяется тестом: $M=\infty$ и $M=$\,маленькое дают байт в байт одинаковый файл.

\section{Параллельная обработка}
\label{sec:parallelism}

Pull выбран в разделе~\ref{sec:approaches}. Здесь --- как этот проход распараллелить, не получив гонок
и не потеряв детерминизм.

\subsection{Без блокировок и детерминизм}
Каждый поток берёт непрерывный диапазон приёмников и пишет только в свой непересекающийся кусок
$s_{\mathrm{new}}$. Массивы $\mathrm{contrib}$ и $s_{\mathrm{old}}$ общие и только на чтение. Отсюда
два следствия.
\begin{itemize}[itemsep=3pt]
  \item \textbf{Горячий цикл без блокировок.} В проходе $\bigO(m)$ нет ни атомиков, ни
        \texttt{synchronized}. Единственная синхронизация --- один барьер на итерацию.
  \item \textbf{Детерминизм, не зависящий от числа потоков.} Никакие два потока не пишут в одну
        ячейку, а порядок сложения внутри ячейки задан порядком рёбер в файле. Поэтому результат
        получается бит в бит одинаковым при $1$ и при $24$ потоках и от запуска к запуску. Планировщик
        потоков на ответ не влияет: ни одна ячейка не делится между потоками, а порядок суммирования в
        ней зафиксирован на диске.
\end{itemize}

\subsection{Балансировка и гипер-узлы}
Делить поровну по числу вершин нельзя. Из-за степенного распределения равные по числу вершин куски
содержат очень разное число рёбер, и поток с гипер-узлом становится отстающим. Поэтому:
\begin{itemize}[nosep]
  \item границы кусков выбираются так, чтобы уравнять число рёбер (по префиксной сумме входящих
        степеней).
  \item кусков создаётся больше, чем потоков (в несколько раз), и пул разбирает их по мере
        освобождения. Это сглаживает остаточный перекос.
  \item гипер-узел в модели pull неделим (у ячейки один владелец), но в абсолютных числах он мал.
        Входящая степень $50\,000$ --- это около $0{,}2$~МБ источников, микросекунды работы.
\end{itemize}

\subsection{Реализация на Java}
\label{sec:java}
\begin{itemize}[itemsep=3pt]
  \item \textbf{Пул потоков и барьер.} Фиксированный \texttt{ExecutorService} на $P$ потоков по числу
        ядер. На каждой итерации задачи-куски отдаются в \texttt{invokeAll}, который ждёт завершения
        всех. Это ожидание и есть барьер итерации: пока все приёмники не пересчитаны, следующая
        итерация не начнётся. Отдельный \texttt{CyclicBarrier} не нужен.
  \item \textbf{Корректность по модели памяти Java.} Потоки пишут в $s_{\mathrm{new}}$ без
        \texttt{volatile} и без блокировок. Это безопасно, потому что завершение задачи и его
        наблюдение через \texttt{invokeAll} задают отношение happens-before. Всё, что поток записал,
        видно главному потоку, когда тот увидел задачу завершённой. Значит записи итерации $t$
        опубликованы до чтений итерации $t{+}1$.
  \item \textbf{Разбиение диапазонами против false sharing.} Поток пишет непрерывный кусок, поэтому
        делиться с соседом может только одна кэш-линия на границе. При кусках в тысячи элементов это
        незаметно. Чередующееся разбиение, наоборот, посадило бы соседние потоки на общие кэш-линии.
  \item \textbf{Детерминированные суммы.} Глобальные суммы (масса фонового узла и норма сходимости)
        считаются в фиксированном порядке: каждый поток отдаёт частичную сумму, а главный поток
        складывает их по порядку номеров. \texttt{DoubleAdder} и \texttt{stream().sum()} с
        неопределённым порядком сознательно не используются. Потоки платформенные, а не виртуальные:
        виртуальные хороши для блокирующего ввода-вывода, но не ускоряют счёт.
  \item \textbf{Мало мусора.} Массивы $s_{\mathrm{old}}$, $s_{\mathrm{new}}$, $\mathrm{contrib}$,
        \texttt{outdeg}, \texttt{sourcesPtr} выделяются один раз. Между итерациями $s_{\mathrm{old}}$
        и $s_{\mathrm{new}}$ просто меняются ссылками. Горячий цикл почти не создаёт мусора, поэтому
        Serial GC справляется даже при крошечной куче.
\end{itemize}

\subsection{Ввод-вывод}
Граф не помещается в память, поэтому файл \texttt{sources} перечитывается с диска на каждой итерации.
Закэшировать его целиком нельзя, в этом суть задачи. Каждый поток читает свои куски последовательно,
держа только буфер фиксированного размера ($64$~КиБ), а не весь кусок. Буфер свой на каждый поток и
переиспользуется между итерациями, поэтому суммарная память на буферы --- это $P\cdot B$, а не $P$
копий чего-то большого. На SSD упреждающее чтение ОС неплохо совмещает ввод-вывод со счётом.

\section{Память и соблюдение лимита}
\label{sec:memory}

\subsection{Принцип out-of-core}
Данные делятся по размеру. Структуры $\bigO(n)$ (по числу вершин) живут в ОЗУ. Структуры $\bigO(m)$
(по числу рёбер) живут на диске и читаются потоком. Поскольку <<гигабайт>> --- это именно рёбра,
потребление памяти определяется числом вершин, а не рёбер. Отсюда следствие, заметное на практике:
размер RSS зависит от числа вершин, а не от размера CSV-файла. Меньший по объёму файл с большим
числом вершин требует больше памяти, чем больший файл с малым числом вершин.

\subsection{Из чего складывается память}
\label{sec:buckets}
Потребление ОЗУ складывается из трёх частей.
\begin{itemize}[itemsep=3pt]
  \item \textbf{Корзина A --- общие векторы $\bigO(n)$.} Три \texttt{double[n]} ($s_{\mathrm{old}}$,
        $s_{\mathrm{new}}$, $\mathrm{contrib}$) и несколько \texttt{int[n]} (\texttt{sourcesPtr},
        \texttt{outdeg}, \texttt{originalIds}) --- это примерно $36$ байт на вершину. Для
        $n\approx 4{,}8$ млн --- около $170$~МБ. Эта часть общая для всех потоков и от их числа не
        зависит.
  \item \textbf{Корзина B --- попоточные буферы.} Один буфер чтения на поток ($64$~КиБ). Всего
        $P\cdot B$ плюс стеки потоков --- единицы МБ. Это единственная часть, растущая с числом
        потоков.
  \item \textbf{Корзина C --- накладные расходы JVM.} Metaspace, кэш кода, структуры сборщика ---
        около $20$--$30$~МБ. Serial GC держит эту часть малой.
\end{itemize}
Основной объём --- корзина A --- общий и не уменьшается с ростом числа потоков. С потоками растёт
только небольшая корзина B.

\subsection{Лимит памяти и его проверка}
\label{sec:cap}
Ограничение памяти задаётся снаружи, а программа подстраивает под него своё поведение. Доступный объём
она узнаёт из максимального размера кучи (значение \texttt{-Xmx}, оно же видно как лимит cgroup в
контейнере) и под него подбирает размер корзин $M$ при предобработке. Один рычаг (\texttt{-Xmx} или
cgroup) задаёт и реальный потолок памяти, и параметры предобработки.

Чтобы симулировать ограничение из условия и проверить, что оно действительно соблюдается, использованы
два механизма.
\begin{itemize}[itemsep=3pt]
  \item \textbf{Лимит кучи JVM (\texttt{-Xmx}).} Ограничивает кучу Java. Нужен, но его одного мало:
        вся память процесса --- это куча плюс стеки, metaspace, кэш кода, буферы вне кучи и страничный
        кэш. Часть из этого лежит вне \texttt{-Xmx}.
  \item \textbf{Лимит всего процесса через cgroup~v2 (основной).} JVM запускается внутри ограничения
        cgroup, и ядро ОС ограничивает всё сразу, снимая процесс по OOM при нарушении:
        \begin{verbatim}
    systemd-run --user --scope -p MemoryMax=256M -p MemorySwapMax=0 \
        java -Xmx128m -XX:+UseSerialGC ...
        \end{verbatim}
        Своп отключён (\texttt{MemorySwapMax=0}), иначе куча незаметно вытеснялась бы в своп и лимит
        стал бы фикцией.
\end{itemize}
Само завершение запуска под лимитом cgroup и есть доказательство: при превышении процесс был бы снят
ядром. Дополнительно в конце запуска печатается пиковый RSS (значение \texttt{VmHWM} из
\texttt{/proc/self/status}) --- получается конкретное число вида <<пик RSS\,$=168$~МиБ при лимите
$256$~МиБ>>.

\paragraph{Деталь про \texttt{-Xmx}.} Он ограничивает только кучу. Поверх неё JVM тратит ещё
$60$--$80$~МБ на нативную часть (сам \texttt{libjvm}, кэш кода, стеки). Поэтому, чтобы удержать RSS
под лимитом $C$, \texttt{-Xmx} стоит брать примерно на $80$~МБ меньше $C$, а не равным $C$.

\subsection{Контрольный контрпример}
Чтобы показать ценность out-of-core, в проект включён намеренно наивный вариант, который держит весь
граф в памяти (вложенные списки в ОЗУ). Под тем же лимитом cgroup он падает, а потоковый движок
завершается. Падает он двумя способами. При куче, равной потоковой, ему не хватает кучи Java
(\texttt{OutOfMemoryError}). При щедрой куче (\texttt{-Xmx1g}) его снимает ядро по OOM (сигнал KILL,
код $137$). Алгоритм и лимит те же, отличается только способ хранения.

\subsection{Ограничения подхода}
\label{sec:limits}
\begin{itemize}[itemsep=3pt]
  \item В ОЗУ обязаны помещаться структуры $\bigO(n)$ --- векторы значений и указатели
        (раздел~\ref{sec:scope}). Если не влезают даже они, это уже другая задача (выгрузка самих
        векторов на диск, как в GraphChi PSW), и она вне рамок работы.
  \item На очень больших графах узкое место --- скорость последовательного чтения с диска (файл
        \texttt{sources} перечитывается на каждой итерации). Компактный формат поднимает этот потолок,
        но саму природу ограничения не убирает.
  \item $128$~МБ --- это тесно для JVM. Нативный пол ($\approx 70$~МБ) плюс векторы $\bigO(n)$ при $n$
        около $2{,}7$ млн ($\approx 108$~МБ) уже превышают $128$~МБ. При таком жёстком лимите это
        аргумент в пользу C++. На Java тот же подход работает при чуть большем лимите.
\end{itemize}

\section{Сложность и точность}
\label{sec:complexity}

Обозначения: $n$ --- число вершин, $m$ --- число рёбер, $P$ --- число потоков, $T$ --- число итераций
до сходимости, $B$ --- размер попоточного буфера.

\subsection{Время}
\paragraph{Предобработка (один раз).} Проход по входу с подсчётом степеней и переименованием ---
$\bigO(m)$. Префиксная сумма для \texttt{sourcesPtr} --- $\bigO(n)$. Раздача рёбер по корзинам ---
$\bigO(m)$. Сортировка корзин --- около $\bigO(m\log(m/K))$ при $K$ корзинах, с малым множителем
(сортируются короткие корзины). Итого предобработка --- $\bigO(m\log m)$ по времени и $\Theta(m)$ по
вводу-выводу (несколько последовательных проходов).

\paragraph{Одна итерация.} Вычисление $\mathrm{contrib}$, инициализация $s_{\mathrm{new}}$, суммы ---
$\bigO(n)$. Ядро~\eqref{eq:pull} --- ровно один проход по файлу \texttt{sources}, $\bigO(m)$ сложений
и $\Theta(m)$ последовательного чтения. Распараллеливание делит работу на $P$ потоков. При
сбалансированных по числу рёбер кусках счётное время итерации --- $\bigO((n+m)/P)$, а время чтения
--- $\Theta(m/\beta)$, где $\beta$ --- эффективная скорость диска. Полное время:
\[
  T_{\text{итер}} = \bigO\!\Big(T\cdot\tfrac{n+m}{P}\Big)\ \text{(счёт)}
    \;+\; \Theta\!\Big(T\cdot\tfrac{m}{\beta}\Big)\ \text{(ввод-вывод)}.
\]
Реальное время --- максимум из двух слагаемых. На больших графах доминирует ввод-вывод ($T$
перечитываний \texttt{sources}), поэтому число читаемых байт сведено к минимуму, а чтение идёт
последовательно.

\subsection{Память}
ОЗУ: $\Theta(n)$ на общие векторы $+\ \bigO(P\cdot B)$ на буферы $+\ \bigO(1)$ накладных JVM. Главное
--- не зависит от $m$ и от наличия гипер-узлов. Диск: $\Theta(m)$ на \texttt{sources} $+\ \Theta(n)$
на остальное.

\subsection{Балансировка: средний и худший случай}
\begin{itemize}[itemsep=2pt]
  \item \textbf{Средний случай.} Куски выровнены по числу рёбер, нагрузка распределена ровно,
        ускорение близко к линейному до потолка ввода-вывода.
  \item \textbf{Худший случай (гипер-узел).} Одна вершина-приёмник степени $D$ неделима, поэтому
        отставание ограничено величиной $\bigO(D)$. При разбиении на $K\gg P$ кусков и $D\ll m/K$ это
        отставание незаметно. Патология наступила бы, только если $D$ сравнима с $m/P$, то есть одна
        вершина собирает заметную долю всех рёбер.
\end{itemize}

\subsection{Сходимость}
Степенной метод на расширенном графе сходится к единственному пределу. Двусторонний фоновый узел
делает цепь Маркова сильно связной и апериодичной, отсюда (теорема Перрона--Фробениуса) существование
и единственность стационарного распределения. Заодно это решает проблему висячих вершин: их масса
уходит фоновому узлу и потом перераспределяется. Ошибка убывает геометрически:
$\|s^{(t)} - s^\ast\|_1 \le C\,\rho^{\,t}$, где $\rho<1$ определяется вторым по модулю собственным
значением матрицы перехода. Число итераций $T = \bigO\!\big(\log(1/\varepsilon)/\log(1/\rho)\big)$.

\subsection{Точность и воспроизводимость}
Все суммы (и в ячейках~\eqref{eq:pull}, и в редукциях) считаются в фиксированном порядке, заданном
раскладкой файла \texttt{sources} и порядком потоков. Сложение \texttt{double} не ассоциативно,
поэтому именно фиксированный порядок даёт побитовую воспроизводимость и независимость от числа
потоков. Накопление ошибки округления контролируется редкой перенормировкой к сохраняемой сумме масс
$n$. Про конкретный критерий остановки и наблюдаемое число итераций на реальных графах --- в
разделе~\ref{sec:testing}.

\section{Тестирование и результаты}
\label{sec:testing}

Тесты решают две задачи: подтвердить корректность и показать работающий результат под лимитом памяти.
Стек --- Gradle, JUnit~5, AssertJ. Всего $106$ тестов, все зелёные.

\subsection{Синтетические тесты (кратко)}
\begin{itemize}[itemsep=3pt]
  \item \textbf{Эталон.} Главный инструмент проверки --- небольшая заведомо верная плотная реализация
        LeaderRank: строится матрица перехода $(n{+}1)\times(n{+}1)$ и считается степенной метод.
        Сложность $\bigO(n^2)$ ограничивает её малыми графами, но на них она даёт <<истину>>, с
        которой сверяется основной движок. На примере \texttt{small.csv} расхождение
        $3{,}3\cdot10^{-15}$, то есть на уровне ошибки округления.
  \item \textbf{Графы вручную.} Треугольник, звезда (это и есть гипер-узел, центр обязан получить
        наибольший ранг), цепочка, висячая вершина, петля, две несвязные компоненты. Каждый случай
        фиксирует своё свойство.
  \item \textbf{Инварианты.} Сохранение массы ($\sum s\approx n$), положительность рангов, сходимость
        за ограниченное число итераций.
  \item \textbf{Главный тест --- инвариантность к числу потоков.} Запуск с $P=1$ и $P=8$ даёт байт в
        байт одинаковый результат. Один этот тест сразу доказывает и корректность параллелизма (нет
        гонок), и детерминизм (нет зависимости от порядка).
  \item \textbf{Генератор данных.} Свой детерминированный генератор R-MAT со степенным распределением
        --- чтобы в графе реально были гипер-узлы. Зерно фиксировано, граф повторяется.
  \item \textbf{Тест под лимитом памяти.} Полный конвейер запускается как отдельный процесс под
        лимитом cgroup. Тест проверяет завершение и пиковый RSS под лимитом. Если на машине нет cgroup
        или \texttt{systemd-run} не может выставить лимит, тест автоматически пропускается, а не
        падает.
\end{itemize}

\subsection{Реальные графы}
Замеры сделаны скриптом \texttt{scripts/benchmark.sh}. Стенд: $24$ ядра, $\approx 29$~ГБ RAM, SSD,
OpenJDK~21, Serial GC. Движок параллельный out-of-core, бюджет памяти берётся из \texttt{-Xmx}. Время
\texttt{compute} --- <<горячее>> (последний прогон при \texttt{-{}-repeat=5}). Время
\texttt{preprocess} однократное, включает разбор CSV и внешнюю сортировку.

\paragraph{Масштабирование по потокам.} Граф R-MAT (scale=20), $16$ млн рёбер, $n = 640\,324$. Число
итераций --- $38$ при любом $P$, результат бит в бит не зависит от $P$.

\begin{center}
\small
\begin{tabular}{rrrrrr}
\toprule
$P$ & preprocess, мс & compute, мс & ускорение & эффективность & пик RSS, МиБ\\
\midrule
1  & 3450{,}9 & 1709{,}7 & 1{,}00 & 100\% & 260{,}6\\
2  & 1911{,}8 & 915{,}5  & 1{,}87 & 93\%  & 220{,}0\\
4  & 1794{,}1 & 543{,}7  & 3{,}14 & 79\%  & 218{,}6\\
8  & 1754{,}6 & 371{,}0  & 4{,}61 & 58\%  & 219{,}7\\
16 & 1758{,}8 & 269{,}2  & 6{,}35 & 40\%  & 222{,}0\\
24 & 1772{,}3 & 231{,}4  & \textbf{7{,}39} & 31\% & 245{,}6\\
\bottomrule
\end{tabular}
\end{center}

Счётное ядро --- это потоковый SpMV, упирающийся в пропускную способность памяти. Поэтому
эффективность падает с ростом $P$: полоса памяти насыщается около $P=8$--$16$. Это честное свойство
задачи, а не баг. Препроцессинг ускоряется примерно вдвое и выходит на плато: остаток --- это
последовательная сборка id-карты и сортировка. Пик RSS почти не зависит от $P$.

\paragraph{Реальные графы SNAP.} $P=24$, \texttt{-Xmx2g}, лимит итераций $100$.

\begin{center}
\small
\begin{tabular}{lrrrrr}
\toprule
граф & $n$ & $m$ & preprocess, мс & compute, мс & пик RSS, МиБ\\
\midrule
web-Google        & $875\,713$   & $5\,105\,039$   & 703{,}9  & 405{,}3  & 237{,}9\\
soc-LiveJournal1  & $4\,847\,571$ & $68\,993\,773$ & 8174{,}7 & 3327{,}9 & 1235{,}1\\
\bottomrule
\end{tabular}
\end{center}

Оба графа считаются out-of-core: рёбра не держатся в памяти. soc-LiveJournal1 с $69$ млн рёбер
обрабатывается под кучей $2$~ГБ.

\paragraph{Сходимость на реальных графах.} Критерий остановки масштабируется числом вершин:
$\sum_i|s_{\mathrm{new}}[i] - s_{\mathrm{old}}[i]| < \mathit{tol}\cdot n$, $\mathit{tol}=10^{-8}$. Это
важно: масса сохраняется и равна $n$, поэтому ненормированная разность растёт с $n$. Фиксированный
абсолютный порог на большом графе требовал бы изменения меньше уровня шума округления и никогда бы не
срабатывал. Деление на $n$ делает <<сошлось>> одинаковым на любом масштабе. Истинное число итераций до
сходимости (замерено поднятием лимита):

\begin{center}
\begin{tabular}{lr}
\toprule
граф & итераций до сходимости\\
\midrule
R-MAT (scale=20)  & 38\\
web-Google        & 191\\
soc-LiveJournal1  & 1909\\
\bottomrule
\end{tabular}
\end{center}

Реальные графы, особенно социальная сеть, имеют меньший спектральный зазор, поэтому степенной метод
сходится медленнее. Оценка <<десятки итераций>> верна для графов вроде R-MAT, но не для соцсети. При
этом ранжирование лидеров стабилизируется намного раньше численной сходимости: top-5 лидеров
совпадает (те же вершины и тот же порядок) при лимите $100$ и при полной сходимости для обоих графов.
Медленный <<хвост>> лишь уточняет младшие разряды очков, не меняя порядок лидеров. Значит лимит $100$
даёт правильных лидеров даже на LiveJournal, а флаг \texttt{-{}-max-iterations} позволяет при желании
довести счёт до полной сходимости.

\paragraph{RSS в зависимости от \texttt{-Xmx}.} На том же R-MAT при сжатии кучи с $2$~ГБ до
$128$~МБ время \texttt{compute} остаётся постоянным ($\approx 215$--$226$~мс), а RSS держится в районе
$190$--$305$~МиБ. Движок не деградирует под давлением памяти, а RSS не следит за размером графа (рёбра
на диске). Главный рычаг RSS --- это \texttt{-Xmx}. Нижняя граница памяти: web-Google считается до
\texttt{-Xmx48m}, soc-LiveJournal1 до \texttt{-Xmx224m}. На границе всё упирается уже в честный
$\bigO(n)$ (векторы значений фазы счёта), а не в предобработку --- ровно тот предел задачи,
обозначенный в разделе~\ref{sec:limits}.

\section{Заключение}
\label{sec:conclusion}

Спроектировано и реализовано вычисление LeaderRank на ориентированном невзвешенном графе, который не
помещается в память. Ключевые решения:
\begin{itemize}[itemsep=2pt]
  \item \textbf{Метрика --- LeaderRank.} Её ядро --- один потоковый проход по рёбрам
        ($A\T\cdot\mathrm{contrib}$), что и делает задачу разрешимой во внешней памяти.
  \item \textbf{Хранение --- CSC по приёмнику.} Указатели \texttt{sourcesPtr} ($\bigO(n)$) в ОЗУ,
        источники \texttt{sources} ($\bigO(m)$) на диске, чтение последовательное.
  \item \textbf{Параллелизм --- модель pull с разбиением по приёмнику.} Без блокировок и
        детерминированно при любом числе потоков.
  \item \textbf{Память --- ограничена числом вершин.} Лимит задаётся снаружи (\texttt{-Xmx}, cgroup),
        программа под него подстраивается, пик RSS измеряется и приводится.
\end{itemize}

Эти решения прямо отвечают критериям задачи: эффективность на больших графах, использование всех
ядер, устойчивость к гипер-узлам, воспроизводимость результата и чистота кода.

Главные числа. Параллельное ядро даёт ускорение $7{,}39\times$ на $24$ потоках, детерминированно
($38$ итераций при любом $P$). Реальные графы SNAP считаются out-of-core: web-Google ($5{,}1$ млн
рёбер) --- около $0{,}4$~с, soc-LiveJournal1 ($69$ млн рёбер) --- около $3{,}3$~с. Контрольный наивный
вариант под тем же лимитом памяти падает, а потоковый движок завершается --- это и подтверждает
ценность out-of-core подхода.


\end{document}
